% Generated draft: requires mathematical and source review.
\documentclass[11pt]{article}
\usepackage[margin=25mm]{geometry}
\usepackage{fontspec}
\setmainfont{Latin Modern Roman}
\usepackage{amsmath,amssymb,mathtools,unicode-math}
\setmathfont{Latin Modern Math}
\usepackage{enumitem,needspace,xurl,hyperref}
\setlength{\parindent}{0pt}
\setlength{\parskip}{5pt}
\setlength{\emergencystretch}{3em}
\newcommand{\N}{\mathbb{N}}
\newcommand{\Z}{\mathbb{Z}}
\newcommand{\Q}{\mathbb{Q}}
\newcommand{\R}{\mathbb{R}}
\newcommand{\C}{\mathbb{C}}
\newcommand{\F}{\mathbb{F}}
\newcommand{\A}{\mathbb{A}}
\newcommand{\PP}{\mathbb{P}}
\newcommand{\E}{\mathbb{E}}
\newcommand{\eps}{\varepsilon}
\newcommand{\dd}{\,\mathrm{d}}
\begin{document}
\section*{UnsolvedMath ID: 11000181}
\textbf{Source problem code:} catalogue entry
\subsection{Definitions and mathematical statement}
A topological Lefschetz pencil (TLP) on an oriented $4$-manifold $X$ is a pencil $X\setminus B\to S^2$ with finite base locus $B$, the usual local pencil model at a base point, and Lefschetz critical points. A regular fiber has a homology class $F\in H_2(X;\Z)$. The exact positivity and equivalence conventions must follow the source.

\textbf{Problem 3.} Given topological descriptions of two terms $f_{k_0}$ and $f_{k_1}$ in the source's asymptotic sequence of pencils on a fixed symplectic $4$-manifold, describe the term
\[
 f_{k_0+k_1}.
\]
The subscripts add; this is not pointwise addition of maps. The excerpt omits the construction of the sequence and its sufficiently-large-index or uniqueness conventions. In the customary setting where the fiber class of $f_k$ is $k\alpha$ for a fixed integral class $\alpha$, the desired output has fiber class $(k_0+k_1)\alpha$. That explanation is conditional on the omitted source convention and does not assert that a homology class uniquely determines a pencil.

The source proposes a broader version: given two pencils on the same symplectic manifold $X$ with fiber classes $F_1,F_2$, describe a TLP with fiber class
\[
 F_1+F_2,
\]
\emph{if one exists}. Existence is not assumed for arbitrary pairs. A useful answer should recover the topological pencil data, for example an appropriate monodromy description under fixed conventions, rather than merely name the sum class. The equivalence relation on input and output descriptions is not stated.

The preceding local item 11000180 motivates this by the Auroux--Katzarkov doubling formula and asks how to access indices beyond repeated doubling. The present source says a sum construction would avoid its specially introduced subclass $\mathcal D^{(0)}_{g,\nu,p}$, whose definition is absent, and might illuminate which classes are represented by symplectic forms on a fixed $4$-manifold.

The appended discussion describes TLP methods for global symplectic topology, mentioning approaches to Taubes-type results independent of Seiberg--Witten theory [6,13,9], extensions to other $4$-manifolds [3], and a possible simple-type application whose precise meaning is unstated. It highlights the canonical class $K(\omega)$ and reports the rational/ruled classification when $\omega\cdot K(\omega)<0$ (equivalently $\theta<0$ in its omitted notation), citing [8]. It then introduces a following request for a pencil-based derivation; that request is not reproduced here. These motivations are not additional solved consequences of the desired sum operation.
\subsection{Short English statement}
Describe a sum operation on topological pencil data that takes two terms of an asymptotic pencil sequence to the term with summed index. More generally, describe a pencil in the sum of two fiber classes whenever such a pencil exists.
\subsection{Sources}
\begin{itemize}
\item Local UnsolvedMath registry, permanent ID 11000181, source code catalogue entry; no proposer or proposal year supplied.
\item Farb, edited volume \emph{Problems on Mapping Class Groups and Related Topics}, chapter ``Lefschetz pencils and mapping class groups,'' Problem 3, PDF page 177; \url{https://www.math.uchicago.edu/~farb/papers/mcgbook.pdf}. The registry records PDF-font-structure extraction and access on 2026-07-29; the inventory distinguishes 329 labeled headings by source page.
\item Adjacent local ID 11000180 gives the asymptotic-sequence and Auroux--Katzarkov doubling motivation. Source labels [6,13,9], [3], and [8] accompany the Taubes, generalization, and rational/ruled discussions; full bibliography entries are absent.
\end{itemize}
\end{document}
